\documentclass{article}
\usepackage[a4paper, margin=2cm]{geometry}
\usepackage{bussproofs}
\usepackage{mathtools}
\usepackage{amssymb}
\usepackage{parskip}
\usepackage{amsthm}


\begin{document}

Given a CCS process $P$ and a finite set of channel names $L \subseteq \mathcal{A}$,
define the operator $\tau_L$ that \emph{hides} every action on a channel in $L$
by replacing it with a silent $\tau$ step. Formally, $\tau_L$ is characterized by:

\[
  \frac{P \xrightarrow{\alpha} P'}{\tau_L(P) \xrightarrow{\tau} \tau_L(P')}
  \quad \alpha \in L \cup \bar{L}
  \qquad
  \frac{P \xrightarrow{\alpha} P'}{\tau_L(P) \xrightarrow{\alpha} \tau_L(P')}
  \quad \alpha \notin L \cup \bar{L}
\]

Show that there exists a CCS context $C_L[\,\cdot\,]$ such that:
\[
\tau_L(P) \sim C_L[P]
\quad \text{for all CCS processes } P
\]
where $\sim$ denotes weak bisimilarity.

The context requires every visible action $a \in L$ to be a $\tau$, using only CCS.
The only source of $\tau$ in CCS is the \textsc{Com} rule: two parallel processes 
synchronize on complementary channels $a$ and $\bar{a}$, producing $\tau$. 

\begin{enumerate}
  \item Rename every $a \in L$ to a fresh private channel $b_a$.
        This gets $L$ out of the way without blocking it.
        \[C_L[P] = P^*\]
        where
        \[P^* = P[\,b_a/a \mid a \in L\,]\]
  \item Run in parallel a partner that is \emph{always} ready to synchronize 
        on $b_a$, in either direction, and for each $a \in L$. 
        \[C_L[P] = P^* \mid A_L\]
        where
        \[A_L = \prod_{a \in L} A_a\]
        and
        \[ A_a \;\overset{\text{def}}{=}\; b_a.A_a + \bar{b}_a.A_a \]
  \item Restrict on $\{b_a\}$ to make $b_a$ invisible from the outside.
        \[ C_L[P] = (P^* \mid A_L)\setminus\{b_a \mid a \in L\} \]
\end{enumerate}

\textbf{Claim.} For all $P \in \text{CCS}$, $\tau_L(P) \sim C_L[P]$.

\begin{proof}
  Two processes are weakly bisimilar if every transition of one can be matched by the other,
  up to silent $\tau$ steps, and the resulting states remain related indefinitely.
  Rather than reasoning about $\sim$ abstractly, we certify this directly: it suffices to
  exhibit a concrete relation $R$ containing the pairs we wish to relate and verify that
  $R$ itself has the mutual-matching property $(*)$ below. Since $\sim$ is the largest
  such relation, any verified $R$ is contained in $\sim$, which gives us our conclusion.

  We therefore take
  \[
    R \;=\; \bigl\{\,(\tau_L(P),\; C_L[P]) \;\bigm|\; P \in \text{CCS}\,\bigr\} \;\cup\; {\sim}
  \]
  The $\cup\;\{\sim\}$ is a technical convenience: after each matching step the residual pair
  need only land back in $R$, and including $\sim$ means any pair already known to be
  bisimilar closes the check immediately without infinite regression.

  Formally, $\sim$ is characterized by: for all processes $P$, $Q$,
  \[
    \begin{array}{rcl}
      P \sim Q &
      \iff &
      \begin{array}{l}
        \text{if } P \xrightarrow{\alpha} P' \;\text{ then }\; \exists\, Q'.\; Q \xRightarrow{\alpha} Q' \;\text{ and }\; P' \sim Q'\\[.5em]
        \text{if } Q \xrightarrow{\alpha} Q' \;\text{ then }\; \exists\, P'.\; P \xRightarrow{\alpha} P' \;\text{ and }\; P' \sim Q'
      \end{array}
    \end{array}
    \tag{$*$}
  \]
  where $\xRightarrow{\alpha}$ denotes weak transition (zero or more $\tau$ steps, 
  then $\alpha$, then zero or more $\tau$ steps; for $\alpha = \tau$ this includes 
  doing nothing).
  
  We now show $R$ satisfies $(*)$. 
  
  \newpage 

  Take any pair $(\tau_L(P),\, C_L[P]) \in R$.

  \begin{itemize}
    \item \textbf{Direction (1): suppose $\tau_L(P) \xrightarrow{\mu} \tau_L(P')$.}\\
          By the defining rules of $\tau_L$ there are two cases:
          \begin{itemize}
            \item \textit{Case $\mu = \tau$, $\alpha \in L$}:
                  Renaming gives $P^* \xrightarrow{b_a} P'^*$; the right branch of $A_a$ provides $\bar{b}_a$:
                  \begin{prooftree}
                    \AxiomC{$P \xrightarrow{a} P'$}
                    \RightLabel{Ren}
                    \UnaryInfC{$P^* \xrightarrow{b_a} P'^*$}
                    \AxiomC{}
                    \RightLabel{Act}
                    \UnaryInfC{$\bar{b}_a.A_a \xrightarrow{\bar{b}_a} A_a$}
                    \RightLabel{Sum}
                    \UnaryInfC{$A_a \xrightarrow{\bar{b}_a} A_a$}
                    \RightLabel{Par-L}
                    \UnaryInfC{$A_L \xrightarrow{\bar{b}_a} A_a \mid A_{L\setminus\{a\}}$}
                    \RightLabel{Com}
                    \BinaryInfC{$P^* \mid A_L \xrightarrow{\tau} P'^* \mid A_L$}
                    \RightLabel{Res}
                    \UnaryInfC{$C_L[P] \xrightarrow{\tau} C_L[P']$}
                  \end{prooftree}
                  So $C_L[P] \xRightarrow{\tau} C_L[P']$ and $(\tau_L(P'), C_L[P']) \in R$.
                  The case $\alpha = \bar{a}$ is symmetric: renaming gives $P^* \xrightarrow{\bar{b}_a} P'^*$
                  and the left branch of $A_a$ provides the complementary $b_a$; the same tree applies.

            \item \textit{Case $\mu = \alpha \notin L \cup \bar{L}$}:
                  Renaming leaves $\alpha$ unchanged and restriction does not block it:
                  \begin{prooftree}
                    \AxiomC{$P \xrightarrow{\alpha} P'$}
                    \RightLabel{Ren}
                    \UnaryInfC{$P^* \xrightarrow{\alpha} P'^*$}
                    \RightLabel{Par-L}
                    \UnaryInfC{$P^* \mid A_L \xrightarrow{\alpha} P'^* \mid A_L$}
                    \RightLabel{Res,\; $\alpha \notin \{b_a \mid a \in L\}$}
                    \UnaryInfC{$C_L[P] \xrightarrow{\alpha} C_L[P']$}
                  \end{prooftree}
                  So $C_L[P] \xRightarrow{\alpha} C_L[P']$ and $(\tau_L(P'), C_L[P']) \in R$.
          \end{itemize}

    \item \textbf{Direction (2): suppose $C_L[P] \xrightarrow{\mu} C_L[P']$.}\\
          Every visible transition must originate from $P^*$.
          $A_L$ alone produces no $\tau$: it has no $\tau$-prefixed subterm and cannot
          self-synchronize, so any $\tau$ arises only by synchronizing with $P^*$ on a
          complementary $b_a$ channel.
          Restriction removes all observable $b_a$-labelled moves from the composed
          system, but does not prevent the internal synchronization that produces $\tau$.
          \begin{itemize}
            \item \textit{Case $\mu = \tau$}:
                  The $\tau$ arises in exactly two ways.

                  \textit{Sub-case: sync on $b_a$ ($P \xrightarrow{a} P'$, $a \in L$):}
                  \begin{prooftree}
                    \AxiomC{$P \xrightarrow{a} P'$}
                    \RightLabel{$a \in L$, first rule of $\tau_L$}
                    \UnaryInfC{$\tau_L(P) \xrightarrow{\tau} \tau_L(P')$}
                  \end{prooftree}
                  The case $P \xrightarrow{\bar{a}} P'$ is symmetric.

                  \textit{Sub-case: silent step ($P \xrightarrow{\tau} P'$):}
                  \begin{prooftree}
                    \AxiomC{$P \xrightarrow{\tau} P'$}
                    \RightLabel{$\tau \notin L \cup \bar{L}$, second rule of $\tau_L$}
                    \UnaryInfC{$\tau_L(P) \xrightarrow{\tau} \tau_L(P')$}
                  \end{prooftree}

                  In both sub-cases $\tau_L(P) \xRightarrow{\tau} \tau_L(P')$
                  and $(\tau_L(P'), C_L[P']) \in R$.

            \item \textit{Case $\mu = \alpha \notin L \cup \bar{L}$}:
                  Then $P \xrightarrow{\alpha} P'$ (renaming leaves $\alpha$ unchanged and
                  restriction does not block it):
                  \begin{prooftree}
                    \AxiomC{$P \xrightarrow{\alpha} P'$}
                    \RightLabel{$\alpha \notin L \cup \bar{L}$, second rule of $\tau_L$}
                    \UnaryInfC{$\tau_L(P) \xrightarrow{\alpha} \tau_L(P')$}
                  \end{prooftree}
                  So $\tau_L(P) \xRightarrow{\alpha} \tau_L(P')$ and $(\tau_L(P'), C_L[P']) \in R$.
          \end{itemize}
  \end{itemize}

In all cases every condition of $(*)$ is satisfied, so $R$ is a weak bisimulation.
Since $\sim$ is the largest weak bisimulation, every bisimulation is contained in $\sim$,
so $R \subseteq \sim$. As $(\tau_L(P), C_L[P]) \in R$, we conclude $\tau_L(P) \sim C_L[P]$ for all $P$.
\end{proof}

\end{document}